
% cf) .................. % !TEX program = xelatex
% Magic comment: 문서 맨 첫 줄에 이 주석을 적어두면, VimTeX(latexmk)이 알아서 xelatex로 컴파일합니다.



% 수식 필수 패키지(amsmath, amssymb)와 한글 패키지(kotex)만 사용하여, pdflatex로 가장 안정적으로 컴파일되는 기본 memoir 템플릿.
% 수식 번호를 붙이고 이를 본문에서 참조(\label과 \ref)하는 핵심 기능만 포함되어 있음.

% :!pdflatex %


\documentclass[12pt,a4paper,oneside]{memoir}


% -- --------------
\usepackage{silence}
% \WarningFilter{latex}{Font shape, LaTeX Font} % 'Font shape'로 시작하는 모든 LaTeX 폰트 경고를 숨깁니다.
% \WarningFilter{latex}{LaTeX Font} % 'Font shape'로 시작하는 모든 LaTeX 폰트 경고를 숨깁니다.
% \WarningFilter{latex}{Font shape} % Font shape로 시작하는 경고 숨기기
\WarningFilter{latex}{LaTeX Font} % Font shape로 시작하는 경고 숨기기
% -- --------------



% --- [최소한의 필수 패키지] ---
\usepackage{kotex}          % 한글 출력을 위한 필수 패키지
\usepackage{amsmath}        % 정렬된 수식(align), 분기(cases) 등 고급 수식용 필수 패키지
\usepackage{amssymb}        % 다양한 수학적 기호 (\mathbb, \partial 등)

% -- --------------

% -- --------------


% -- --------------

% --- [memoir 기본 레이아웃 설정] ---
% 기본 여백을 상하좌우 3cm로 깔끔하게 조절합니다.
\setlrmarginsandblock{3cm}{3cm}{*} % 좌우 여백
\setulmarginsandblock{3cm}{3cm}{*} % 上하 여백
\checkandfixthelayout

% 장(Chapter) 스타일을 너무 거대하지 않게, 일반 섹션 스타일처럼 단정하게 바꿉니다.
\chapterstyle{section}

% --- [문서 정보 설정] ---
\title{vimtex test in nvim}
\author{summer\_willow}
\date{\today}

\begin{document}

% 제목 출력
\maketitle

% 목차 출력 (컴파일을 2번 연속 해야 번호가 올바르게 갱신됩니다)
\tableofcontents
\newpage

\chapter{vimtex 기능}

VimTeX의 기본 접두사(LocalLeader)는 거꾸로 된 슬래시인 $\backslash$ (백슬래시)입니다. 


\begin{itemize}
\item \verb|\ll| (VimtexCompile): 실시간 컴파일 시작/중지 \\
(Background에서 latexmk가 상시 구동됨)
처음 누르면 배경에서 컴파일이 시작되고, 이후 파일을 저장(:w)할 때마다 자동으로 컴파일되어 PDF가 갱신됩니다. 다시 누르면 컴파일러가 꺼집니다.

\item \verb|\lv| (VimtexView): PDF 뷰어 열기 및 순방향 검색 \\
컴파일된 PDF를 Sumatra PDF로 띄우고, 현재 커서가 위치한 곳으로 화면을 이동시킵니다.

\item \verb|\le| (VimtexErrors): 에러 로그 확인 \\
문법 에러 등으로 컴파일이 실패했을 때 에러 창을 (Quickfix) 열어 원인을 보여줍니다. 토글방식

\item \verb|\lc| (VimtexClean): 임시 파일 삭제 \\
컴파일 과정에서 생기는 .aux, .log, .synctex.gz 등의 자잘한 임시 파일들을 깔끔하게 지워줍니다.

\end{itemize}


\section{Inverse sync}

3. SumatraPDF 설정 (역방향 검색: PDF $\rightarrow$ Neovim 이동)

PDF를 더블 클릭했을 때 Neovim의 해당 코드 라인으로 바로 이동하게 만드는 설정

SumatraPDF > 설정(Settings) > 옵션(Options) > '역방향 검색 명령줄 설정 (Set inverse search command-line)' 

(※ 만약 nvim-qt나 특정 터미널 에뮬레이터를 사용하여 항상 하나의 Neovim 인스턴스로만 띄우고 싶다면 neovim-remote 파이썬 패키지를 설치한 뒤 nvr 명령어를 사용하기도 하지만, 위 headless 방식이 윈도우 환경에서 설정 오류가 가장 적고 깔끔합니다).


\section{매직 코멘트(Magic Comment)}
\verb|% !TEX program = xelatex|

VimTeX는 기본적으로 latexmk라는 프로그램에게 컴파일을 맡깁니다. 이 latexmk는 아무런 지시가 없으면 기본값인 pdflatex 엔진으로 빌드를 시도합니다.

문제점: pdflatex는 한글(kotex 패키지 등)을 컴파일할 때 에러가 나거나 한글이 깨집니다.

해결책: 그렇기 때문에 문서 맨 위에 \verb|% !TEX program = xelatex|를 적어서 "이 문서는 pdflatex 말고 xelatex로 구워줘!"라고 명시하는 것입니다.

\subsection{매직 코멘트 없이 항상 xelatex로 빌드하는 방법 (추천)}

매번 파일 상단에 매직 코멘트를 적는 게 귀찮으시다면, Neovim 설정에서 latexmk가 항상 xelatex를 쓰도록 강제하면 됩니다.

plugins/vimtex.lua 파일 설정을 아래와 같이 수정


% \begin{verbatim}
% vim.g.vimtex_compiler_latexmk = {
%       options = {
%         "-xelatex", -- 항상 xelatex 엔진을 사용하도록 설정
%         "-file-line-error",
%         "-synctex=1",
%         "-interaction=nonstopmode",
%       },
%     }
% \end{verbatim}


111 222 333
aaa

bbb

ddd

ccc

eee
333
444
555








\chapter{해석학 및 신호 처리 기초 수식}
% \chapter{해석학 및 신호 처리 기초 수式}

\section{수식 정렬과 번호 부여}
LaTeX에서 가장 기본적이면서도 중요한 수식 환경은 \texttt{align}입니다. 등호(\texttt{=}) 앞에 \texttt{\&} 기호를 붙여 위치를 맞출 수 있으며, 각 줄마다 자동으로 수식 번호가 부여됩니다.

푸리에 변환(Fourier Transform)의 기본 정의와 선형성 성질은 다음과 같이 나타낼 수 있습니다.
\begin{align}
    X(f) &= \int_{-\infty}^{\infty} x(t) e^{-j 2\pi f t} dt \label{eq:fourier} \\
    \mathcal{F}\{\alpha x(t) + \beta y(t)\} &= \alpha X(f) + \beta Y(f) \label{eq:linearity}
\end{align}

본문에서 특정 수식을 언급하고 싶을 때는 \texttt{\textbackslash ref\{\}}를 사용합니다. 예를 들어, 식 \ref{eq:fourier}는 연속 시간 신호의 주파수 영역 해석을 위한 기본 정의식이며, 식 \ref{eq:linearity}는 그에 따른 선형성 성질을 나타냅니다. 이렇게 적어두면 나중에 수식이 추가되어 번호가 밀리더라도 LaTeX가 알아서 번호를 찾아줍니다.


\section{번호를 생략하는 수식과 상자 구조}
만약 본문 중간에 강조하고 싶지만 굳이 수식 번호를 붙여 참조할 필요가 없는 경우에는 \texttt{equation*} 환경이나 \texttt{\textbackslash[ ... \textbackslash]}를 사용합니다.

오일러 공식(Euler's formula)은 다음과 같습니다.
\[
    e^{j\theta} = \cos\theta + j\sin\theta
\]

조건에 따라 값이 나뉘는 함수를 표현할 때는 \texttt{cases} 환경을 사용합니다. 이 역시 하나의 완성된 수식이므로 식 \ref{eq:step_func}과 같이 끝에 라벨을 붙여 참조할 수 있습니다.
\begin{equation}
    u(t) = \begin{cases}
        1 & \text{if } t \ge 0 \\
        0 & \text{if } t < 0
    \end{cases} \label{eq:step_func}
\end{equation}

식 \ref{eq:step_func}은 신호 처리와 제어 공학에서 가장 기본적으로 사용되는 단위 계단 함수(Unit Step Function)의 정의입니다.


\chapter{고급 수식 및 대형 연산자}

\section{행렬과 벡터 표현}
다차원 데이터를 다룰 때 필수적인 행렬은 \texttt{bmatrix} 환경을 사용해 각 괄호\texttt{[]} 형태로 표현합니다.
\begin{equation}
    \begin{bmatrix}
        y_1 \\
        y_2
    \end{bmatrix} = 
    \begin{bmatrix}
        a & b \\
        c & d
    \end{bmatrix}
    \begin{bmatrix}
        x_1 \\
        x_2
    \end{bmatrix} + 
    \begin{bmatrix}
        n_1 \\
        n_2
    \end{bmatrix} \label{eq:matrix_system}
\end{equation}
식 \ref{eq:matrix_system}과 같은 선형 연립방정식 체계는 디지털 통신 채널 모델링의 가장 기초가 되는 형태입니다.

\section{다중 적분과 극한 수식}
미적분학에서 자주 등장하는 대형 연산자(인테그랄, 시그마 등)와 아래위 첨자 배치 예시입니다.
\begin{equation}
    \lim_{n \to \infty} \sum_{k=1}^{n} \frac{1}{n} f\left(\frac{k}{n}\right) = \int_{0}^{1} f(x) dx
\end{equation}

\end{document}
